
An \emph{overlap-compatible modification record} below retains the
actual analytic charts, all centres and blowup maps, and the
isomorphisms used to reglue them. In a local replacement
$X=U\cup V$, the modified chart $U'$ has a specified overlap
$W'\simeq W=U\cap V$; this isomorphism is the restriction of
the recorded composite map on $U$. The replacement is
$X'=U'\cup_W V$. Thus this datum specifies an analytic
operation, including its neighbourhoods. It is not merely a
list of component graphs or normal degrees. The records used
here are the finite local $\sigma$-records in
\cite[\S\S4.2--4.14]{Kulikov1977}, with the corrections and
extra fibre blowups in \cite[\S\S8--11]{Kulikov1981}.

\begin{proposition}[The Hopf-minimum branch is generated by canonical-flat cores and exact sigma records]\label{cdp:st:C.2.UN.75}
Let $f:X\to\Delta$ be a proper degeneration of connected
complex two-tori, with smooth total threefold and full
reduced simple-normal-crossing central fibre $S=\sum_iV_i$.
Choose a generator $\Omega$ of $f_*K_X$ near zero and write
\[
             \operatorname{div}(\Omega)=\sum_i k_iV_i,
                  \qquad k_i\ge0,\quad \min_i k_i=0.
\tag{C.KG.1}\label{cdp:eq:C.KG.1}
\]
Let $\mathcal H$ be the set of indices whose component has
a Hopf surface as its smooth minimal model. The following
two conditions are equivalent:
\begin{enumerate}
\item $k_i=0$ for every $i\in\mathcal H$.
\item There is a finite overlap-compatible modification
record from $f$ to a smooth proper torus degeneration
$f_0:X_0\to\Delta$ whose full central fibre $S_0$ is
reduced simple NC and whose canonical line is trivial.
\end{enumerate}
The constructed comparison is an actual bimeromorphic map
over the same disc and induces an isomorphism
\[
                X\setminus S\ \simeq\ X_0\setminus S_0
                         \quad\hbox{over }\Delta^*.
\tag{C.KG.2}\label{cdp:eq:C.KG.2}
\]
Intermediate charts retain their normalization branches;
in particular a type-II move may temporarily give a
self-intersection of a global component.

The endpoint has precisely one of the following forms,
with the actual markings and line maps of the indicated
statements:
\begin{enumerate}
\item a smooth two-torus;
\item a cycle of the marked diagonal Hopf and degree-zero
elliptic ruled pairs in \ref{cdp:st:C.2.UN.65}, with dual edge normals and
the actual complement-lattice product $M=I$;
\item a positive toric central log space in \ref{cdp:st:C.2.UN.71}, given by
$(N,\Gamma,\mathscr P,c)$, with Euler number
$2[N:\Gamma]$ and the specified full toric component pairs.
\end{enumerate}
Every degeneration in the stated branch is recovered by
reversing its finite record on its actual endpoint
neighbourhood. Conversely every such reverse record with
a smooth full reduced simple-NC final fibre gives a
degeneration in this branch. These are inverse analytic
constructions, retaining every centre, overlap and normal
extension.

Here is the exact Euler bookkeeping for these records.
Orient a recorded blowup arrow positively from its target
to its blowup, and negatively in the reverse direction.
For its actual centre $L$ and ideal $I$, put
\[
 \begin{split}
  E_I&=\operatorname{Proj}_L\left(\bigoplus_{a\ge0}
                                      I^a/I^{a+1}\right),\\
  \delta(I)&=e_c(E_I)-e_c(L)
       =\int_L\bigl(e((E_I)_x)-1\bigr)\,de_c(x).
 \end{split}
\tag{C.KG.3}\label{cdp:eq:C.KG.3}
\]
Local terms are evaluated on a finite constructible
subdivision subordinate to the retained charts. Restored
overlaps cancel. If $\epsilon_\nu\in\{1,-1\}$ records
the directions in the forward reduction to $S_0$, then
\[
             e(S)=e(S_0)-\sum_\nu\epsilon_\nu\delta(I_\nu).
\tag{C.KG.4}\label{cdp:eq:C.KG.4}
\]
For an ordinary smooth centre in a smooth threefold,
$\delta=2$ for a point and $\delta=2-2g(L)$ for a compact
curve. Both type-I and type-II moves have total defect
zero. An ordinary ruled-section contraction in \ref{cdp:st:C.2.UN.74} has
Euler change $-(2-2g(B))$ and removes its specified number
of triple points. Thus the exact zero-Euler criterion is
\[
 \sum_\nu\epsilon_\nu\delta(I_\nu)=
 \begin{cases}
 0,&S_0\text{ is of type 1 or 2},\\
 2[N:\Gamma],&S_0\text{ is of type 3}.
 \end{cases}
\tag{C.KG.5}\label{cdp:eq:C.KG.5}
\]
The normal-cone formula also applies to auxiliary
singular-centre charts; their contribution is not
replaced by a smooth-curve formula.

For a given compact total space over a compact curve,
the reduction and its inverse take place on that same
compact base, by the actual isomorphism \eqref{cdp:eq:C.KG.2} on an
annulus. A reverse record on an actual compact core
realization of \ref{cdp:st:C.2.UN.65} or \ref{cdp:st:C.2.UN.73} likewise gives a compact
realization. Its total Euler change is the same as
the central-fibre change in \eqref{cdp:eq:C.KG.4}.
\end{proposition}

\begin{proof}
\par\medskip\noindent\textbf{The canonical generator and its normalized orders.}
The proper direct image $f_*K_X$ is a torsion-free
rank-one coherent sheaf on the disc, hence a line.
On $\Delta^*$ its restriction and base change are the
one-dimensional spaces of top forms on the tori.
A generator therefore has no zero away from $S$.
Its divisor is the effective divisor in \eqref{cdp:eq:C.KG.1}.
If every $k_i$ were positive, $\Omega/f$ would be
holomorphic on the smooth total space, since the
\emph{full} fibre is reduced. This would contradict
that $\Omega$ generates the direct-image line.
Consequently its minimum is zero. Another generator
is $a(t)\Omega$ for a holomorphic unit $a$, so these
orders are independent of that choice.

For later comparison, let $\phi:X\dashrightarrow X'$
be an actual bimeromorphic map over the disc between
smooth total spaces. Pullback of a holomorphic top
form is holomorphic wherever $\phi$ is holomorphic.
The indeterminacy has codimension at least two;
Hartogs extension of the coefficient in smooth
coordinate volume frames gives a holomorphic form
on all of $X$. The inverse map gives the inverse
pullback, since their compositions agree on their
common biholomorphic open set. The construction
works over every open subdisc, and gives the actual
isomorphism
\[
                f'_*K_{X'}\xrightarrow{\ \phi^*\ } f_*K_X.
\tag{C.KG.6}\label{cdp:eq:C.KG.6}
\]
In particular normalized generators agree under
pullback up to a base unit. At the generic point
of a divisor which survives the comparison,
the map is a biholomorphism of the corresponding
divisorial local rings, so its canonical order is
unchanged.

\par\medskip\noindent\textbf{Why every Hopf-birational component is retained.}
A holomorphic map from $\PP^1$ to a Hopf surface
is constant: lift it to the universal cover
$\C^2\setminus\{0\}$ and use the two constant
coordinate functions. Every rational curve on
a point blowup of that surface is therefore
exceptional over the Hopf minimal model. A finite
sequence of point blowups has only finitely many
such curves. These surfaces cannot be covered by
rational curves. They are also not projective:
their first Betti number is one, unchanged by
point blowups, whereas a projective surface has
even first Betti number.

In the recorded local $\sigma$-maps, a divisorial
exceptional component over a general point of a
curve centre is ruled by complete rational
fibres. A component lying over an isolated
point is projective, as a component of a
projectivized normal cone. Here the contradiction
with the original compact component uses the
proper composite graphs, not a comparison of
local analytic patches. Suppose a strict
transform were removed for the first time.
For a curve centre, choose a general complete
rational exceptional fibre through a generic
strict-transform point. The proper preceding
graph maps this curve meromorphically back to
the original compact component. The map is
nonconstant, since at that generic point the
strict-transform comparison is a local
biholomorphism. Varying the fibre gives
rational curves through a dense open subset
of the original component, contradicting the
finite exceptional-curve description above.
For a point centre, restrict the proper
preceding graph to the compact projective
exceptional component containing the generic
strict-transform piece. Its image back in
the original compact component is closed
and contains that open piece, and therefore
is the whole component. The comparison has
generic degree one. Resolving this surface
graph gives a bimeromorphic comparison with
a projective surface; point blowups preserve
$b_1$, again contradicting $b_1=1$.
Thus the strict transforms through the
finite chart record retain a prime divisor
of that birational type. The restored
overlap maps identify these strict transforms
by their actual generic-point maps. Hence a
component in $\mathcal H$ survives to the
endpoint.

At a canonical-flat smooth endpoint a nowhere-zero
three-form generates $f_{0*}K_{X_0}$, so every
canonical order is zero. Equation \eqref{cdp:eq:C.KG.6} then
proves $k_i=0$ for every $i\in\mathcal H$.
This proves the necessary minimum condition for
the full Hopf birational class, including
arbitrary point blowups. No identification of
different conventions for the notation
\emph{VII$_0$} is used here.

\par\medskip\noindent\textbf{The finite reduction and the actual allowed charts.}
\ref{cdp:st:C.2.UN.19} identifies every actual class-VII minimal
model in a torus degeneration as Hopf. Thus the
minimum condition just stated implies the
condition in \cite[Theorem~1']{Kulikov1981}.
That theorem applies to proper analytic good
degenerations with trivial fibre canonical
class; its hypotheses do not include
projectivity. The definition of good in
\cite[p.958]{Kulikov1977} is exactly the smooth
total space and full reduced NC fibre with
smooth components used here.

We specify the comparison maps in that proof.
For a nonboundary $(-1)$-curve $A\subset V_i$
meeting one double curve transversely once
and avoiding the remaining boundary, the
principal-fibre normal sequence is
\[
 0\longrightarrow\OO_{\PP^1}(-1)
 \longrightarrow N_{A/X}
 \longrightarrow\OO_{\PP^1}(-1)
 \longrightarrow0.
\tag{C.KG.7}\label{cdp:eq:C.KG.7}
\]
It splits because $H^1(\PP^1,\OO)=0$.
Blow up $A$; the exceptional divisor is
$\PP^1\times\PP^1$ with normal
$\OO(-1,-1)$. Contract its other ruling
using the corrected contraction theorem in
\ref{cdp:st:C.2.UN.64}. This transfers the point blowup to
the neighbour. The exact adjunction equation
gives $k_i=k_j$ for these two components.
This is the actual type-I map.

For a double curve which is $(-1)$ on both
sides, the same normal bundle calculation
gives the same two-ruling construction.
The principal-fibre line has exactly two
triple points on that curve. If the opposite
components there have orders $k_3,k_4$,
adjunction gives
\[
                         k_1+k_2=k_3+k_4.
\tag{C.KG.8}\label{cdp:eq:C.KG.8}
\]
The ordinary blowup has exceptional full
multiplicity two; the other contraction
removes the old double curve and introduces
the new one between the two opposite
branches. If those branches belong to one
global component, both normalization branches
are retained, as in
\cite[\S4.4]{Kulikov1977}. One does not
discard that self-incidence or call this
intermediate fibre embedded simple NC.

The other contraction charts are the actual
ruled section-and-fibre maps of \ref{cdp:st:C.2.UN.74} and
the plane-line map of \ref{cdp:st:C.2.UN.64}. For a reducible
ruling fibre, the local contractions and
inverse-order restorations in
\cite[\S\S4.10.7--4.10.9,4.14.5]{Kulikov1977}
give the retained finite chart record.
The corrected proof in
\cite[\S8]{Kulikov1981} also blows up the
specified missing fibre components before
contracting the chosen ruled surface.
The new ruled components have canonical
order $k+1$ when the removed component has
order $k$. For a maximum-order component
no such additions are needed.

Here termination is an actual finite
induction, rather than a decrease of an
unspecified complexity. Let $R=\max k_i$.
Every removed old component has level
$k\in\{1,\ldots,R\}$. It produces only
finitely many new components, all at level
$k+1$, and produces none when $k=R$.
There are finitely many level-one removals;
therefore finitely many level-two objects
are ever produced. Continuing through the
fixed finite set of levels proves that all
removals terminate. The components left
have order zero. The corrected endpoint
argument in \cite[\S10]{Kulikov1981}, using
\cite[\S4.17]{Kulikov1977}, gives a smooth
total space and smooth reduced components.
In particular \eqref{cdp:eq:C.KG.8} prevents a minimum
component from acquiring a self-incidence
in a type-II move whose removed edge has
an order above the minimum. The endpoint
generator is therefore nowhere zero.

\par\medskip\noindent\textbf{The source's actual overlap maps glue global graphs.}
The bare two-cover definition of modification
in \cite[\S2]{Kulikov1981} is not used to
infer a global map. The construction itself
provides the following additional data.
In \cite[\S4.10.8, pp.980--982]{Kulikov1977}
one keeps $V$ fixed, performs the displayed
finite contractions and restorations in
$U$, and obtains $U'$ with an overlap
$W'$ isomorphic to the original $W$ by
the restriction of that same composite.
Section4.14.5 repeats this construction.
The added local fibre blowups of
\cite[\S8]{Kulikov1981} leave this overlap
unchanged.

For a finite $\sigma$-chain, its composite
graph is obtained as the closure of the
common biholomorphic graph in the proper
image of the fibre product of its arrow
graphs. Both projections are proper.
Over $W$ this graph is precisely the
graph of the restored isomorphism.
Glue it to the identity graph over $V$.
The graphs agree on their overlap, and
properness is local on each projection
target. The glued graph is consequently
proper over both $U\cup V$ and
$U'\cup_W V$. It gives actual mutually
inverse bimeromorphic maps. All its arrows
are over the same disc and all their
centres are central. Their restriction
over $\Delta^*$ is an isomorphism.
Iteration proves \eqref{cdp:eq:C.KG.2} and justifies
the generator comparison \eqref{cdp:eq:C.KG.6}.

This argument does not assert a single
global factorization into ordinary
blowups. The finite open covers and
their induced overlap maps are part
of the stated modification record.
Auxiliary singular charts in the
restoration are also retained rather
than silently replaced by smooth ones.

\par\medskip\noindent\textbf{Endpoint classification and the converse record.}
Apply \ref{cdp:st:C.2.UN.65} to the actual endpoint.
Its zero-Euler alternatives give types1
and2, including all the marked line
and integral-lattice comparisons.
Its positive case is \ref{cdp:st:C.2.UN.68}; \ref{cdp:st:C.2.UN.71}
identifies its actual central log
space with the specified toric data
by the displayed chart and residue
maps. This proves the stated complete
endpoint list.

Reverse each recorded arrow and each
overlap isomorphism. They reconstruct
the original neighbourhood literally.
Starting instead with any actual
canonical-flat endpoint and such a
reverse record, all off-central maps
remain isomorphisms; the general
fibres are the same connected tori.
At a smooth full reduced final fibre
the argument following \eqref{cdp:eq:C.KG.6} shows
that every Hopf-birational component
has order zero. This proves both
directions of the generation statement.
In this statement the endpoint is an
actual analytic family. The maps do
not replace its neighbourhood by a
different realization merely because
their central log spaces agree.

\par\medskip\noindent\textbf{The exact Euler defect, including singular centres.}
For an actual blowup with centre ideal
$I$, its exceptional inverse image is
the projectivized normal cone in
\eqref{cdp:eq:C.KG.3}. Outside $L$ it is an
isomorphism. Since $L$ is central,
the reduced central inverse image
replaces precisely $L$ by $E_I$.
Additivity of compactly supported
Euler characteristic proves \eqref{cdp:eq:C.KG.3}
for this arrow, including its
constructible-fibre integral.
For these computations refine the
retained analytic charts by relatively
compact subanalytic coordinate
neighbourhoods, with closures inside
the original charts. Compactness of
the central fibre gives a finite
subcover. Subdivide its overlaps and
centres compatibly, and use inverse
images under the retained proper
graphs; all pieces have finite
subanalytic triangulations. This
refinement also defines the local
terms when a displayed centre is an
open part of a curve. Euler numbers
of arbitrary analytic open covers
are not used.

In one replacement, evaluated on this
finite-type refinement,
$U\cup V\rightsquigarrow
U'\cup_WV$, inclusion--exclusion gives
\[
 \begin{split}
 e(S')-e(S)
   &=e_c(S'\cap U')-e_c(S\cap U)\\
   &=\sum_{\nu\text{ in this chart}}
                         \epsilon_\nu\delta(I_\nu),
 \end{split}
\tag{C.KG.9}\label{cdp:eq:C.KG.9}
\]
because the unchanged $V$ and the
isomorphic restored overlaps cancel.
Equivalently one applies this equality
to the finite subdivision just
described and sums; no boundary
piece is counted twice. The local
arrow contributions telescope, and
then the global replacements
telescope, proving \eqref{cdp:eq:C.KG.4}.

For a smooth centre of dimension $a$
in a smooth threefold the normal
cone is its rank-$(3-a)$ normal
bundle, so its exceptional fibre
is $\PP^{2-a}$. Thus
\[
                \delta(I)=(2-a)e_c(L).
\tag{C.KG.10}\label{cdp:eq:C.KG.10}
\]
For typeI and typeII the common
blowup replaces a compact $\PP^1$
by $\PP^1\times\PP^1$ from either
side, so the two contributions
$+2$ and $-2$ cancel. The ruled
and plane changes follow from
\ref{cdp:st:C.2.UN.74} and \ref{cdp:st:C.2.UN.64}. In an auxiliary
singular-centre chart \eqref{cdp:eq:C.KG.3}
remains the formula: for example,
the blowup of the three coordinate
axes in $\C^3$, with ideal
$(xy,xz,yz)$, has fibre $\PP^1$
over a general point of each axis
and fibre $\PP^2$ over their
common origin. The last defect
is two rather than one. This
illustrates why the full normal
cone was retained.
Substitution of the three endpoint
Euler numbers proves \eqref{cdp:eq:C.KG.5}.

\par\medskip\noindent\textbf{Compact scope.}
Restrict the given compact family
to a small disc at the chosen
value. On a smaller annulus use
the actual isomorphism \eqref{cdp:eq:C.KG.2}
to glue the new proper filling
to the unchanged complement.
The two proper pieces give a
Hausdorff proper analytic family;
their smooth charts coincide on
the overlap. This produces the
claimed compact replacement,
and the inverse record recovers
the original compact family.
The same construction starts
from the explicit compact cores
of \ref{cdp:st:C.2.UN.65} and \ref{cdp:st:C.2.UN.73} and any reverse
record on their actual tubes.
Euler additivity cancels the
unchanged complement and annulus,
so the total-space Euler change
is precisely the central one.
\end{proof}


